%Preamble
\documentclass[12pt]{article}
\title{Puncte Foarte Generale de Urmărit ale Proiectului}
\author{Dumitru Marius-Vlad}
\date{\today}

% Document
\begin{document}
\maketitle

Principalele puncte foarte generare ale proiectului sunt:
    \begin{itemize}
        \item Dezvoltarea unuia sau mai multor agenți, la alegere. Fie unul mare care sa le facă pe toate, fie mai mulți agenți micuți care să facă 1 singur task simplu(probabil alegerea potrivită este mai mulți agenți micuți care să facă 1 singur task.)
        \item Folosirea tehnicilor și/sau algoritmilor de Reinforcement Learning clasic și/sau Deep Reinforcement Learning, la alegere.
        \item Folosirea unui singur joc care va reprezenta mediul în care agenții vor interacționa. Agenți vor fi participanți(jucători) la meciuri iar interacțiunile vor fi de tip jucător$\rightarrow$jucător și jucător$\rightarrow$mediu(joc).
        \item Agentul/Agenții vor fi jucători, prin urmare, el/ei trebuie să fie capabili să perceapă situația curentă, să decidă ce și cum (în baza situației curente), să joace(să execute acțiuni) pentru a atinge un scop(țel) final.
        \item În cele din urmă, menirea agentului/agenților este obținerea(atingerea) unui scop(țel) final. Prin urmare, tot procesul de percepție $\rightarrow$ decizie $\rightarrow$ acțiune trebuie să urmărească acest scop(țel) final.
        \item Scopul(țelul) final trebuie foarte bine definit, foarte clar. Deci lanțul percepție $\rightarrow$ decizie $\rightarrow$ acțiune va fi într-o direcție pentru obținerea unui anumit scop și într-o altă direcție diferită pentru obținerea altui scop diferit.
        \item Mai multe experimente ce acoperă scenarii diferite. Scopurile agenților trebuie să acopere multiple situații/scenarii distincte. Acest aspect dă o diversitate proiectului.
        \item Proiectul trebuie să explice foarte bine lanțul percepție $\rightarrow$ decizie $\rightarrow$ acțiune. DE CE anume sa decis și actionat într-un anume fel, în anume condiții vs. decis și acționat altfel în alte condiții. Răspuns la întrebarea generică "DE CE, ÎN SITUAȚIA X, SA PROCEDAT ÎN ACEST MOD ȘI NU ALTFEL / DE CE, ÎN SITUAȚIA X, SA DECIS ACEST MOD DE OPERARE SI NU ALTFEL ?" (Explicabilitate)
        \item Antrenarea agenților cu ajutorul seturilor de date relevante. Crearea de noi seturi de date(daca este cazul) / extinderea celor existente (daca este cazul).
        \item Evaluarea agenților prin metrici specifice RL/DRL. OBLIGATORIU trebuiesc avute în vedere metricile timpi de decizie, rată de succes (trecute în cerințele lucrării). Multiple experimente.
        \item Comparații atât între agenții dezvoltați în proiect cât și cu literatura de specialitate. Eventual propuneri de îmbunătățire.
        \item Proiectul trebuie să fie funcțional și ușor de folosit.
        \item ( Opțional ). Folosirea de teste pentru verificarea codului (Test Driven Development).
        \item Folosire predominant Python + pachete relevante.
        \item Dacă este nevoie, dar într-o măsură limitată, folosirea și altor limbaje de programare. Dominant trebuie să rămână python + pachete.
        \item Agentul/Agenții și jocul sunt entități diferite. Jocul trebuie să ofere agentului/agenților, informații despre situația curentă a meciului și posibilitatea de manipulare a jucătorilor. Deci, între joc și agent/agenți trebuie să existe o interfață pentru interacțiune. Dacă jocul nu o oferă nativ, aceasta trebuie construită. (integrarea agenților cu mediul)
    \end{itemize}
\end{document}